\documentclass[12pt]{article}
\usepackage[utf8]{inputenc}
\usepackage[T1]{fontenc}
\usepackage{hyperref}
\title{Sou Nama Language Resources for K-12 Curriculum Development}
\author{Compiled by Claude Sonnet 5 (claude-sonnet-5) for Andr\'as Kornai}
\date{}
\begin{document}
\maketitle

\section{Introduction}
Sou Nama (ISO 639-3 tlt), also called Teluti, is an Austronesian language of south Seram, Maluku, Indonesia. Search results link it with the Sepa--Teluti entry and show some confusion between Teluti, Sepa and the unrelated Sou Upaa (Manusela). Speaker figures range from about 4,600 (ASJP) to 18,000--20,000 (SIL-based sources and Wikipedia). I found no detailed grammar or school materials.

\section{Primary Educational Resources}
None found.

\section{Linguistic and Reference Materials}
\begin{itemize}
\item Wikipedia, Sepa--Teluti language: \url{https://en.wikipedia.org/wiki/Sepa%E2%80%93Teluti_language}
\item Numerals wordlist for Sou Nama (Max Planck EVA): \url{https://lingweb.eva.mpg.de/channumerals/Sou-Nama.htm}
\item Glottolog reference, Stokhof 1981 (Holle wordlists): \url{https://glottolog.org/resource/reference/id/148049}
\end{itemize}

\section{Religious and Community Resources}
None found.

\section{Digital and Audio Resources}
\begin{itemize}
\item ASJP wordlist, Taluti: \url{https://asjp.clld.org/languages/TALUTI}
\end{itemize}

\section{Curriculum Development Context}
Only catalogue entries and wordlists turned up. Any curriculum would need new documentation, and the Teluti/Sepa identity should be clarified first.

\end{document}